% Figure: the relative thickness of the velocity and thermal boundary layers
% for three Prandtl-number regimes. For Pr ~ 1 (gases) the two layers grow
% together; for Pr >> 1 (oils) the thermal layer is buried inside a thick
% velocity layer; for Pr << 1 (liquid metals) the thermal layer outruns the
% velocity layer. The ratio delta/delta_t ~ Pr^(1/3) is the source of the
% Prandtl power in every forced-convection correlation.
\begin{figure}[htbp]
\centering
\begin{tikzpicture}
\begin{axis}[
  width=13cm, height=5.2cm,
  xlabel={distance along plate $x$ (arb.)},
  ylabel={layer thickness (arb.)},
  xmin=0, xmax=10, ymin=0, ymax=3.2,
  axis lines=left,
  legend pos=north west,
  legend cell align=left,
  tick label style={font=\scriptsize},
  label style={font=\small},
  legend style={font=\scriptsize},
  domain=0:10, samples=60,
]
% velocity layer (common reference, ~ sqrt(x))
\addplot[engblue, very thick] {0.55*sqrt(x)};
\addlegendentry{velocity layer $\delta$}
% thermal layer, Pr ~ 1 (air): comparable
\addplot[englime, very thick, dashed] {0.60*sqrt(x)};
\addlegendentry{thermal $\delta_t$, $\text{Pr}\approx 1$ (air)}
% thermal layer, Pr >> 1 (oil): thinner
\addplot[engred, very thick, dotted] {0.24*sqrt(x)};
\addlegendentry{thermal $\delta_t$, $\text{Pr}\gg 1$ (oil)}
% thermal layer, Pr << 1 (liquid metal): thicker
\addplot[engorange, very thick, dash dot] {1.05*sqrt(x)};
\addlegendentry{thermal $\delta_t$, $\text{Pr}\ll 1$ (liquid metal)}
\end{axis}
\end{tikzpicture}
\caption[Velocity and thermal boundary-layer thickness by Prandtl number]{The
velocity boundary layer (the region of slowed fluid) and the thermal boundary
layer (the region of changed temperature) grow together along a flat plate,
each thickening as the square root of distance. The Prandtl number fixes their
ratio, $\delta/\delta_t \sim \text{Pr}^{1/3}$. A gas with $\text{Pr}\approx 1$
carries the two layers at comparable thickness; a high-Prandtl oil buries a
thin thermal layer deep inside a thick velocity layer, so the wall temperature
gradient is steep and the convection coefficient large for the same velocity
field; a liquid metal with $\text{Pr}\ll 1$ has a thermal layer far thicker than
the velocity layer because heat diffuses faster than momentum. The fractional
power of $\text{Pr}$ in every flat-plate Nusselt correlation is this thickness
ratio in disguise.}
\label{fig:vol04ch05:prandtl-layers}
\end{figure}
